% Синтетический кейс: hyperref + cleveref (\cref, \Cref, \autoref, \hyperref), метки с @cref в .aux.
\documentclass[12pt]{article}
\usepackage[T2A]{fontenc}
\usepackage[utf8]{inputenc}
\usepackage[russian]{babel}
\usepackage{amsmath,amsthm}
\usepackage[unicode]{hyperref}
\usepackage[russian]{cleveref}

\newtheorem{theorem}{Теорема}[section]
\crefname{theorem}{теорема}{теоремы}
\Crefname{theorem}{Теорема}{Теоремы}
\crefname{equation}{формула}{формулы}
\Crefname{equation}{Формула}{Формулы}

\begin{document}

\section{Ссылки}\label{sec:refs}

\begin{theorem}\label{thm:c}
Теорема для ссылок.
\end{theorem}

\begin{equation}\label{eq:c}
  y = z .
\end{equation}

Обычные ссылки: теорема~\ref{thm:c}, формула~\eqref{eq:c}, раздел~\ref{sec:refs}.
Умные ссылки: \cref{thm:c}, \Cref{eq:c}, \cref{sec:refs}; \autoref{thm:c};
\hyperref[thm:c]{текст ссылки}; \url{https://example.org}.

\begin{theorem}\label{thm:d}
Вторая теорема.
\end{theorem}
\begin{gather}
  a = 1, \label{eq:d}\\
  b = 2. \label{eq:e}
\end{gather}

Несколько: \cref{thm:c,thm:d}, \cref{eq:c,eq:d,eq:e}; \cref{thm:c,eq:c}.
Диапазон: \crefrange{eq:c}{eq:e}.

\end{document}
